\section{Discussion}
\label{sec:discussion}

\subsection*{Key Findings}

\textbf{Finding 1: The graph coordinate system matters more than the symmetry group
for cross-architecture SSL.}
Graph-based SSL with permutation equivariance achieves ${\sim}220\%$ $R^2$ improvement over
flat tokenization.
The representational gap is mechanistic ---
SANE collapse vs.\ EquiSSL-perm's distributed latent space.
Equally important: scale equivariance does not improve over permutation-only
($\Delta R^2 = -0.143$ in the h-m2 ablation evaluation;
$\Delta R^2 = -0.021$ in the h-m1 evaluation under comparable training conditions).
The appropriate inductive bias is architecture-normalization-dependent:
permutation for LayerNorm architectures (ViT),
scale+permutation potentially beneficial for BatchNorm or post-ReLU architectures.

\textbf{Finding 2: MMD is an unreliable metric for cross-architecture transfer evaluation.}
SANE achieves lower MMD than graph-based encoders despite substantially worse $R^2$.
Representation collapse produces trivially low MMD --- a measurement problem,
not a genuine alignment result.
More robust alternatives include per-dimension KL divergence, Fr\'{e}chet distance with
covariance estimation, or $R^2$-based linear probe directly.

\textbf{Finding 3: Property prediction and functional model generation are separable
capabilities.}
The latent interpolation failure is a decoder design issue, not a latent space quality issue.
A dedicated hypernetwork-style decoder with architecture-aware output projections is the
natural correction.

\subsection*{Limitations}

\textbf{Statistical power.}
All cross-architecture $R^2$ results are from seed 0 only, with 53 ViT-S/16 models.
With $n = 1$ seed, statistical significance is not assessable.
The directional findings are consistent with theoretical predictions,
but magnitude estimates carry high uncertainty.
Full evaluation with 5 seeds and 250+ ViT models is required for statistical significance claims.

\textbf{Ablation fairness.}
The h-m2 ablation (Table~\ref{tab:ablation}) compares EquiSSL from h-e1 (50 epochs)
against EquiSSL-perm from h-m1 (100 epochs).
This training duration discrepancy is a confound:
the observed $\Delta R^2 = +0.143$ advantage for EquiSSL-perm may partly reflect additional
training rather than the symmetry group difference.
The h-m1 evaluation (Table~\ref{tab:main}) provides a more controlled comparison
and shows a smaller but directionally consistent advantage ($\Delta R^2 = +0.021$).

\textbf{Scale equivariance reversal.}
The LayerNorm gauge-fixing hypothesis is empirically motivated but not directly tested as a
causal mechanism.
Direct validation requires evaluating scale vs.\ permutation-only on non-LayerNorm
architectures (e.g., ReLU MLP, BatchNorm CNN) as target.

\textbf{Training data.}
We use approximately 3{,}000 of ${\sim}30{,}000$ available SANE MultiZoo CNN checkpoints.
Full dataset training may improve $R^2$.

\textbf{Decoder design.}
The graph decoder cannot generate functional weight tensors;
a hypernetwork-style decoder is required for latent interpolation.

\textbf{MMD comparability.}
Cross-architecture MMD (h-e1) was measured for SANE and EquiSSL only.
No cross-architecture MMD is available for EquiSSL-perm,
limiting the distribution shift analysis (RQ2) to two methods.

\subsection*{Broader Impact}

Weight-space SSL methods that generalize across architecture families could reduce the cost
of model zoo analysis --- enabling property prediction for arbitrary pre-trained models
without dedicated training data for each new architecture.
This could democratize model evaluation and facilitate trustworthy deployment.
Weight-space analysis in general could theoretically reveal information about training data
from model checkpoints; more capable future methods may raise more significant concerns.
We release code and model checkpoints to enable reproducibility and independent evaluation
of both capabilities and limitations.
